% paper85-agent-planning-descent.tex — Agent Planning as Descent:
%   Verifying Multi-Step Plans Before Execution
%   Paper 85 of the Judgment Geometry series.
% Compile: pdflatex paper85-agent-planning-descent
\documentclass[11pt]{article}
\usepackage{lmodern}
% jugeo-common.tex — shared preamble for all Judgment Geometry papers
% Include via: % jugeo-common.tex — shared preamble for all Judgment Geometry papers
% Include via: % jugeo-common.tex — shared preamble for all Judgment Geometry papers
% Include via: \input{jugeo-common}

% ── Packages ────────────────────────────────────────────────────────────
\usepackage{amsmath,amssymb,amsthm,mathtools}
\usepackage{tikz-cd}
\usepackage{listings}
\usepackage{xcolor}
\usepackage{xspace}
\usepackage{booktabs}
\usepackage{multirow}
\usepackage{enumitem}
\usepackage[expansion=false]{microtype}
\usepackage{hyperref}
\usepackage{cleveref}
% \usepackage{thmtools}  % disabled: not available in this TeX installation
\usepackage{float}
\usepackage{caption}
\usepackage{subcaption}
\usepackage{tabularx}
\usepackage{stmaryrd}       % \llbracket, \rrbracket
\usepackage{bbm}            % \mathbbm{1}

% ── Page geometry ───────────────────────────────────────────────────────
\usepackage[margin=1in]{geometry}

% ── Theorem environments ────────────────────────────────────────────────
\theoremstyle{plain}
\newtheorem{theorem}{Theorem}[section]
\newtheorem{lemma}[theorem]{Lemma}
\newtheorem{proposition}[theorem]{Proposition}
\newtheorem{corollary}[theorem]{Corollary}

\theoremstyle{definition}
\newtheorem{definition}[theorem]{Definition}
\newtheorem{example}[theorem]{Example}
\newtheorem{construction}[theorem]{Construction}

\theoremstyle{remark}
\newtheorem{remark}[theorem]{Remark}
\newtheorem{notation}[theorem]{Notation}

% ── Python listings style ──────────────────────────────────────────────
\definecolor{jugeoBlue}{HTML}{2563EB}
\definecolor{jugeoGreen}{HTML}{059669}
\definecolor{jugeoRed}{HTML}{DC2626}
\definecolor{jugeoPurple}{HTML}{7C3AED}
\definecolor{jugeoGray}{HTML}{6B7280}
\definecolor{jugeoBg}{HTML}{F8FAFC}

\lstdefinestyle{jugeo-python}{
  language=Python,
  basicstyle=\ttfamily\small,
  keywordstyle=\color{jugeoBlue}\bfseries,
  stringstyle=\color{jugeoGreen},
  commentstyle=\color{jugeoGray}\itshape,
  numberstyle=\tiny\color{jugeoGray},
  backgroundcolor=\color{jugeoBg},
  numbers=left,
  numbersep=6pt,
  frame=single,
  framerule=0.4pt,
  rulecolor=\color{jugeoGray!40},
  breaklines=true,
  breakatwhitespace=true,
  showstringspaces=false,
  tabsize=4,
  xleftmargin=1.5em,
  framexleftmargin=1.5em,
  morekeywords={dataclass,frozen,field,Enum,IntEnum,auto,Any,
                Mapping,Sequence,Optional,match,case},
  literate={→}{{$\to$}}1 {←}{{$\leftarrow$}}1
           {≤}{{$\leq$}}1 {≥}{{$\geq$}}1 {≠}{{$\neq$}}1
           {∈}{{$\in$}}1 {∀}{{$\forall$}}1 {∃}{{$\exists$}}1
           {⊕}{{$\oplus$}}1 {⊖}{{$\ominus$}}1,
  escapeinside={(*@}{@*)},
}
\lstset{style=jugeo-python}

\lstdefinestyle{jugeo-output}{
  basicstyle=\ttfamily\small,
  backgroundcolor=\color{jugeoBg},
  frame=single,
  framerule=0.4pt,
  rulecolor=\color{jugeoGray!40},
  breaklines=true,
  showstringspaces=false,
  xleftmargin=1.5em,
  framexleftmargin=0.5em,
  numbers=none,
}

% ── JuGeo-specific macros ──────────────────────────────────────────────

% Judgment 8-tuple
\newcommand{\J}{\mathcal{J}}
\newcommand{\Jud}[8]{(#1,\,#2,\,#3,\,#4,\,#5,\,#6,\,#7,\,#8)}
\newcommand{\JudShort}[1]{\J_{#1}}

% Components
\newcommand{\coord}{\mathsf{c}}            % coordinate
\newcommand{\prop}{\varphi}                 % proposition
\newcommand{\carrier}{\mathcal{A}}          % carrier type
\newcommand{\evid}{\mathcal{E}}             % evidence bundle
\newcommand{\oblig}{\mathcal{O}}            % obligations
\newcommand{\obstr}{\mathcal{B}}            % obstructions
\newcommand{\trust}{\mathcal{T}}            % trust annotation
\newcommand{\prov}{\Pi}                     % provenance

% Trust algebra
\newcommand{\Talg}{\mathfrak{T}}            % trust ordered algebra
\newcommand{\Eadm}{\mathcal{E}_{\mathrm{adm}}}
\newcommand{\tleq}{\preceq}                 % trust ordering
\newcommand{\tjoin}{\oplus}                 % conservative join
\newcommand{\tatten}{\ominus}               % attenuation
\newcommand{\tpromote}{\uparrow_\pi}        % promotion
\newcommand{\tdemote}{\downarrow_\chi}       % demotion

% Trust tiers
\newcommand{\Tcontra}{\ensuremath{\mathsf{CONTRADICTED}}}
\newcommand{\Tunver}{\ensuremath{\mathsf{UNVERIFIED}}}
\newcommand{\Tcopilot}{\ensuremath{\mathsf{COPILOT}}}
\newcommand{\Toracle}{\ensuremath{\mathsf{ORACLE}}}
\newcommand{\Truntime}{\ensuremath{\mathsf{RUNTIME}}}
\newcommand{\Tsolver}{\ensuremath{\mathsf{SOLVER}}}
\newcommand{\Tproof}{\ensuremath{\mathsf{PROOF}}}

% Site and geometry
\newcommand{\Site}{\mathcal{S}}             % semantic site
\newcommand{\Cov}{\mathrm{Cov}}            % covering family
\newcommand{\Ob}{\mathrm{Ob}}              % objects
\newcommand{\Mor}{\mathrm{Mor}}            % morphisms
\newcommand{\res}{\mathrm{res}}            % restriction
\newcommand{\Cech}{\check{C}}              % Čech complex
\newcommand{\Hcoh}[1]{H^{#1}}             % cohomology group

% Descent
\newcommand{\descent}{\mathrm{desc}}
\newcommand{\glue}{\mathrm{glue}}
\newcommand{\overlap}[2]{#1 \cap #2}

% Semantic moves
\newcommand{\VERIFY}{\textsc{Verify}}
\newcommand{\CONSTRUCT}{\textsc{Construct}}
\newcommand{\NORMALIZE}{\textsc{Normalize}}
\newcommand{\GLUE}{\textsc{Glue}}
\newcommand{\RESTRICT}{\textsc{Restrict}}
\newcommand{\TRANSPORT}{\textsc{Transport}}
\newcommand{\REPAIR}{\textsc{Repair}}
\newcommand{\PROMOTE}{\textsc{Promote}}
\newcommand{\CHALLENGE}{\textsc{Challenge}}
\newcommand{\ESCALATE}{\textsc{Escalate}}

% Fragments
\newcommand{\QFLIA}{\texttt{QF\_LIA}}
\newcommand{\QFLRA}{\texttt{QF\_LRA}}
\newcommand{\QFBV}{\texttt{QF\_BV}}
\newcommand{\QFUF}{\texttt{QF\_UF}}

% Misc
\newcommand{\Python}{\textsc{Python}}
\newcommand{\JuGeo}{\textsc{JuGeo}}
\newcommand{\LEAN}{\textsc{Lean}\,4}
\newcommand{\Fstar}{F$^{\star}$}
\newcommand{\Coq}{\textsc{Coq}}
\newcommand{\Dafny}{\textsc{Dafny}}

% Common shortcuts
\newcommand{\ie}{i.e.\@\xspace}
\newcommand{\eg}{e.g.\@\xspace}
\newcommand{\cf}{cf.\@\xspace}
\newcommand{\wrt}{w.r.t.\@\xspace}

% Evaluation shortcuts
\newcommand{\numcases}{300}
\newcommand{\numspec}{100}
\newcommand{\numeq}{100}
\newcommand{\numbug}{100}
\newcommand{\medtime}{6.5\,ms}
\newcommand{\totaltime}{1.949\,s}


% ── Packages ────────────────────────────────────────────────────────────
\usepackage{amsmath,amssymb,amsthm,mathtools}
\usepackage{tikz-cd}
\usepackage{listings}
\usepackage{xcolor}
\usepackage{xspace}
\usepackage{booktabs}
\usepackage{multirow}
\usepackage{enumitem}
\usepackage[expansion=false]{microtype}
\usepackage{hyperref}
\usepackage{cleveref}
% \usepackage{thmtools}  % disabled: not available in this TeX installation
\usepackage{float}
\usepackage{caption}
\usepackage{subcaption}
\usepackage{tabularx}
\usepackage{stmaryrd}       % \llbracket, \rrbracket
\usepackage{bbm}            % \mathbbm{1}

% ── Page geometry ───────────────────────────────────────────────────────
\usepackage[margin=1in]{geometry}

% ── Theorem environments ────────────────────────────────────────────────
\theoremstyle{plain}
\newtheorem{theorem}{Theorem}[section]
\newtheorem{lemma}[theorem]{Lemma}
\newtheorem{proposition}[theorem]{Proposition}
\newtheorem{corollary}[theorem]{Corollary}

\theoremstyle{definition}
\newtheorem{definition}[theorem]{Definition}
\newtheorem{example}[theorem]{Example}
\newtheorem{construction}[theorem]{Construction}

\theoremstyle{remark}
\newtheorem{remark}[theorem]{Remark}
\newtheorem{notation}[theorem]{Notation}

% ── Python listings style ──────────────────────────────────────────────
\definecolor{jugeoBlue}{HTML}{2563EB}
\definecolor{jugeoGreen}{HTML}{059669}
\definecolor{jugeoRed}{HTML}{DC2626}
\definecolor{jugeoPurple}{HTML}{7C3AED}
\definecolor{jugeoGray}{HTML}{6B7280}
\definecolor{jugeoBg}{HTML}{F8FAFC}

\lstdefinestyle{jugeo-python}{
  language=Python,
  basicstyle=\ttfamily\small,
  keywordstyle=\color{jugeoBlue}\bfseries,
  stringstyle=\color{jugeoGreen},
  commentstyle=\color{jugeoGray}\itshape,
  numberstyle=\tiny\color{jugeoGray},
  backgroundcolor=\color{jugeoBg},
  numbers=left,
  numbersep=6pt,
  frame=single,
  framerule=0.4pt,
  rulecolor=\color{jugeoGray!40},
  breaklines=true,
  breakatwhitespace=true,
  showstringspaces=false,
  tabsize=4,
  xleftmargin=1.5em,
  framexleftmargin=1.5em,
  morekeywords={dataclass,frozen,field,Enum,IntEnum,auto,Any,
                Mapping,Sequence,Optional,match,case},
  literate={→}{{$\to$}}1 {←}{{$\leftarrow$}}1
           {≤}{{$\leq$}}1 {≥}{{$\geq$}}1 {≠}{{$\neq$}}1
           {∈}{{$\in$}}1 {∀}{{$\forall$}}1 {∃}{{$\exists$}}1
           {⊕}{{$\oplus$}}1 {⊖}{{$\ominus$}}1,
  escapeinside={(*@}{@*)},
}
\lstset{style=jugeo-python}

\lstdefinestyle{jugeo-output}{
  basicstyle=\ttfamily\small,
  backgroundcolor=\color{jugeoBg},
  frame=single,
  framerule=0.4pt,
  rulecolor=\color{jugeoGray!40},
  breaklines=true,
  showstringspaces=false,
  xleftmargin=1.5em,
  framexleftmargin=0.5em,
  numbers=none,
}

% ── JuGeo-specific macros ──────────────────────────────────────────────

% Judgment 8-tuple
\newcommand{\J}{\mathcal{J}}
\newcommand{\Jud}[8]{(#1,\,#2,\,#3,\,#4,\,#5,\,#6,\,#7,\,#8)}
\newcommand{\JudShort}[1]{\J_{#1}}

% Components
\newcommand{\coord}{\mathsf{c}}            % coordinate
\newcommand{\prop}{\varphi}                 % proposition
\newcommand{\carrier}{\mathcal{A}}          % carrier type
\newcommand{\evid}{\mathcal{E}}             % evidence bundle
\newcommand{\oblig}{\mathcal{O}}            % obligations
\newcommand{\obstr}{\mathcal{B}}            % obstructions
\newcommand{\trust}{\mathcal{T}}            % trust annotation
\newcommand{\prov}{\Pi}                     % provenance

% Trust algebra
\newcommand{\Talg}{\mathfrak{T}}            % trust ordered algebra
\newcommand{\Eadm}{\mathcal{E}_{\mathrm{adm}}}
\newcommand{\tleq}{\preceq}                 % trust ordering
\newcommand{\tjoin}{\oplus}                 % conservative join
\newcommand{\tatten}{\ominus}               % attenuation
\newcommand{\tpromote}{\uparrow_\pi}        % promotion
\newcommand{\tdemote}{\downarrow_\chi}       % demotion

% Trust tiers
\newcommand{\Tcontra}{\ensuremath{\mathsf{CONTRADICTED}}}
\newcommand{\Tunver}{\ensuremath{\mathsf{UNVERIFIED}}}
\newcommand{\Tcopilot}{\ensuremath{\mathsf{COPILOT}}}
\newcommand{\Toracle}{\ensuremath{\mathsf{ORACLE}}}
\newcommand{\Truntime}{\ensuremath{\mathsf{RUNTIME}}}
\newcommand{\Tsolver}{\ensuremath{\mathsf{SOLVER}}}
\newcommand{\Tproof}{\ensuremath{\mathsf{PROOF}}}

% Site and geometry
\newcommand{\Site}{\mathcal{S}}             % semantic site
\newcommand{\Cov}{\mathrm{Cov}}            % covering family
\newcommand{\Ob}{\mathrm{Ob}}              % objects
\newcommand{\Mor}{\mathrm{Mor}}            % morphisms
\newcommand{\res}{\mathrm{res}}            % restriction
\newcommand{\Cech}{\check{C}}              % Čech complex
\newcommand{\Hcoh}[1]{H^{#1}}             % cohomology group

% Descent
\newcommand{\descent}{\mathrm{desc}}
\newcommand{\glue}{\mathrm{glue}}
\newcommand{\overlap}[2]{#1 \cap #2}

% Semantic moves
\newcommand{\VERIFY}{\textsc{Verify}}
\newcommand{\CONSTRUCT}{\textsc{Construct}}
\newcommand{\NORMALIZE}{\textsc{Normalize}}
\newcommand{\GLUE}{\textsc{Glue}}
\newcommand{\RESTRICT}{\textsc{Restrict}}
\newcommand{\TRANSPORT}{\textsc{Transport}}
\newcommand{\REPAIR}{\textsc{Repair}}
\newcommand{\PROMOTE}{\textsc{Promote}}
\newcommand{\CHALLENGE}{\textsc{Challenge}}
\newcommand{\ESCALATE}{\textsc{Escalate}}

% Fragments
\newcommand{\QFLIA}{\texttt{QF\_LIA}}
\newcommand{\QFLRA}{\texttt{QF\_LRA}}
\newcommand{\QFBV}{\texttt{QF\_BV}}
\newcommand{\QFUF}{\texttt{QF\_UF}}

% Misc
\newcommand{\Python}{\textsc{Python}}
\newcommand{\JuGeo}{\textsc{JuGeo}}
\newcommand{\LEAN}{\textsc{Lean}\,4}
\newcommand{\Fstar}{F$^{\star}$}
\newcommand{\Coq}{\textsc{Coq}}
\newcommand{\Dafny}{\textsc{Dafny}}

% Common shortcuts
\newcommand{\ie}{i.e.\@\xspace}
\newcommand{\eg}{e.g.\@\xspace}
\newcommand{\cf}{cf.\@\xspace}
\newcommand{\wrt}{w.r.t.\@\xspace}

% Evaluation shortcuts
\newcommand{\numcases}{300}
\newcommand{\numspec}{100}
\newcommand{\numeq}{100}
\newcommand{\numbug}{100}
\newcommand{\medtime}{6.5\,ms}
\newcommand{\totaltime}{1.949\,s}


% ── Packages ────────────────────────────────────────────────────────────
\usepackage{amsmath,amssymb,amsthm,mathtools}
\usepackage{tikz-cd}
\usepackage{listings}
\usepackage{xcolor}
\usepackage{xspace}
\usepackage{booktabs}
\usepackage{multirow}
\usepackage{enumitem}
\usepackage[expansion=false]{microtype}
\usepackage{hyperref}
\usepackage{cleveref}
% \usepackage{thmtools}  % disabled: not available in this TeX installation
\usepackage{float}
\usepackage{caption}
\usepackage{subcaption}
\usepackage{tabularx}
\usepackage{stmaryrd}       % \llbracket, \rrbracket
\usepackage{bbm}            % \mathbbm{1}

% ── Page geometry ───────────────────────────────────────────────────────
\usepackage[margin=1in]{geometry}

% ── Theorem environments ────────────────────────────────────────────────
\theoremstyle{plain}
\newtheorem{theorem}{Theorem}[section]
\newtheorem{lemma}[theorem]{Lemma}
\newtheorem{proposition}[theorem]{Proposition}
\newtheorem{corollary}[theorem]{Corollary}

\theoremstyle{definition}
\newtheorem{definition}[theorem]{Definition}
\newtheorem{example}[theorem]{Example}
\newtheorem{construction}[theorem]{Construction}

\theoremstyle{remark}
\newtheorem{remark}[theorem]{Remark}
\newtheorem{notation}[theorem]{Notation}

% ── Python listings style ──────────────────────────────────────────────
\definecolor{jugeoBlue}{HTML}{2563EB}
\definecolor{jugeoGreen}{HTML}{059669}
\definecolor{jugeoRed}{HTML}{DC2626}
\definecolor{jugeoPurple}{HTML}{7C3AED}
\definecolor{jugeoGray}{HTML}{6B7280}
\definecolor{jugeoBg}{HTML}{F8FAFC}

\lstdefinestyle{jugeo-python}{
  language=Python,
  basicstyle=\ttfamily\small,
  keywordstyle=\color{jugeoBlue}\bfseries,
  stringstyle=\color{jugeoGreen},
  commentstyle=\color{jugeoGray}\itshape,
  numberstyle=\tiny\color{jugeoGray},
  backgroundcolor=\color{jugeoBg},
  numbers=left,
  numbersep=6pt,
  frame=single,
  framerule=0.4pt,
  rulecolor=\color{jugeoGray!40},
  breaklines=true,
  breakatwhitespace=true,
  showstringspaces=false,
  tabsize=4,
  xleftmargin=1.5em,
  framexleftmargin=1.5em,
  morekeywords={dataclass,frozen,field,Enum,IntEnum,auto,Any,
                Mapping,Sequence,Optional,match,case},
  literate={→}{{$\to$}}1 {←}{{$\leftarrow$}}1
           {≤}{{$\leq$}}1 {≥}{{$\geq$}}1 {≠}{{$\neq$}}1
           {∈}{{$\in$}}1 {∀}{{$\forall$}}1 {∃}{{$\exists$}}1
           {⊕}{{$\oplus$}}1 {⊖}{{$\ominus$}}1,
  escapeinside={(*@}{@*)},
}
\lstset{style=jugeo-python}

\lstdefinestyle{jugeo-output}{
  basicstyle=\ttfamily\small,
  backgroundcolor=\color{jugeoBg},
  frame=single,
  framerule=0.4pt,
  rulecolor=\color{jugeoGray!40},
  breaklines=true,
  showstringspaces=false,
  xleftmargin=1.5em,
  framexleftmargin=0.5em,
  numbers=none,
}

% ── JuGeo-specific macros ──────────────────────────────────────────────

% Judgment 8-tuple
\newcommand{\J}{\mathcal{J}}
\newcommand{\Jud}[8]{(#1,\,#2,\,#3,\,#4,\,#5,\,#6,\,#7,\,#8)}
\newcommand{\JudShort}[1]{\J_{#1}}

% Components
\newcommand{\coord}{\mathsf{c}}            % coordinate
\newcommand{\prop}{\varphi}                 % proposition
\newcommand{\carrier}{\mathcal{A}}          % carrier type
\newcommand{\evid}{\mathcal{E}}             % evidence bundle
\newcommand{\oblig}{\mathcal{O}}            % obligations
\newcommand{\obstr}{\mathcal{B}}            % obstructions
\newcommand{\trust}{\mathcal{T}}            % trust annotation
\newcommand{\prov}{\Pi}                     % provenance

% Trust algebra
\newcommand{\Talg}{\mathfrak{T}}            % trust ordered algebra
\newcommand{\Eadm}{\mathcal{E}_{\mathrm{adm}}}
\newcommand{\tleq}{\preceq}                 % trust ordering
\newcommand{\tjoin}{\oplus}                 % conservative join
\newcommand{\tatten}{\ominus}               % attenuation
\newcommand{\tpromote}{\uparrow_\pi}        % promotion
\newcommand{\tdemote}{\downarrow_\chi}       % demotion

% Trust tiers
\newcommand{\Tcontra}{\ensuremath{\mathsf{CONTRADICTED}}}
\newcommand{\Tunver}{\ensuremath{\mathsf{UNVERIFIED}}}
\newcommand{\Tcopilot}{\ensuremath{\mathsf{COPILOT}}}
\newcommand{\Toracle}{\ensuremath{\mathsf{ORACLE}}}
\newcommand{\Truntime}{\ensuremath{\mathsf{RUNTIME}}}
\newcommand{\Tsolver}{\ensuremath{\mathsf{SOLVER}}}
\newcommand{\Tproof}{\ensuremath{\mathsf{PROOF}}}

% Site and geometry
\newcommand{\Site}{\mathcal{S}}             % semantic site
\newcommand{\Cov}{\mathrm{Cov}}            % covering family
\newcommand{\Ob}{\mathrm{Ob}}              % objects
\newcommand{\Mor}{\mathrm{Mor}}            % morphisms
\newcommand{\res}{\mathrm{res}}            % restriction
\newcommand{\Cech}{\check{C}}              % Čech complex
\newcommand{\Hcoh}[1]{H^{#1}}             % cohomology group

% Descent
\newcommand{\descent}{\mathrm{desc}}
\newcommand{\glue}{\mathrm{glue}}
\newcommand{\overlap}[2]{#1 \cap #2}

% Semantic moves
\newcommand{\VERIFY}{\textsc{Verify}}
\newcommand{\CONSTRUCT}{\textsc{Construct}}
\newcommand{\NORMALIZE}{\textsc{Normalize}}
\newcommand{\GLUE}{\textsc{Glue}}
\newcommand{\RESTRICT}{\textsc{Restrict}}
\newcommand{\TRANSPORT}{\textsc{Transport}}
\newcommand{\REPAIR}{\textsc{Repair}}
\newcommand{\PROMOTE}{\textsc{Promote}}
\newcommand{\CHALLENGE}{\textsc{Challenge}}
\newcommand{\ESCALATE}{\textsc{Escalate}}

% Fragments
\newcommand{\QFLIA}{\texttt{QF\_LIA}}
\newcommand{\QFLRA}{\texttt{QF\_LRA}}
\newcommand{\QFBV}{\texttt{QF\_BV}}
\newcommand{\QFUF}{\texttt{QF\_UF}}

% Misc
\newcommand{\Python}{\textsc{Python}}
\newcommand{\JuGeo}{\textsc{JuGeo}}
\newcommand{\LEAN}{\textsc{Lean}\,4}
\newcommand{\Fstar}{F$^{\star}$}
\newcommand{\Coq}{\textsc{Coq}}
\newcommand{\Dafny}{\textsc{Dafny}}

% Common shortcuts
\newcommand{\ie}{i.e.\@\xspace}
\newcommand{\eg}{e.g.\@\xspace}
\newcommand{\cf}{cf.\@\xspace}
\newcommand{\wrt}{w.r.t.\@\xspace}

% Evaluation shortcuts
\newcommand{\numcases}{300}
\newcommand{\numspec}{100}
\newcommand{\numeq}{100}
\newcommand{\numbug}{100}
\newcommand{\medtime}{6.5\,ms}
\newcommand{\totaltime}{1.949\,s}

\input{data-paper85}

% ── Additional packages ────────────────────────────────────────────
\usepackage{array}
\usepackage{tikz}
\usetikzlibrary{arrows.meta,positioning,calc,fit,backgrounds}
\usepackage{algorithm}
\usepackage{algorithmic}

% ── Paper-specific macros ──────────────────────────────────────────
\newcommand{\PlanSheaf}{\mathcal{P}}
\newcommand{\GoalSite}{\Site_{\mathsf{G}}}
\newcommand{\SubGoal}[1]{U_{#1}}
\newcommand{\PlanVerify}{\textsc{PlanVerify}}
\newcommand{\PlanRepair}{\textsc{PlanRepair}}
\newcommand{\ActionSeq}{\mathcal{A}}
\newcommand{\PreCond}{\mathrm{Pre}}
\newcommand{\PostCond}{\mathrm{Post}}
\newcommand{\GoalSpace}{\mathcal{G}}
\newcommand{\CoverFam}{\{U_i\}_{i \in I}}
\newcommand{\CocycleCond}{\mathrm{Cocycle}}
\newcommand{\ObsClass}{\mathrm{Obs}}
\newcommand{\RepairOp}{\mathcal{R}}
\newcommand{\ConflictSet}{\mathcal{C}}
\newcommand{\PlanCost}{\mathrm{cost}}
\newcommand{\Severity}{\mathrm{sev}}
\newcommand{\PlanGlue}{\mathrm{glue}}
\newcommand{\PlanRes}{\mathrm{res}}
\newcommand{\DescentData}{\mathrm{DD}}
\newcommand{\GoalCat}{\mathbf{Goal}}
\newcommand{\ActCat}{\mathbf{Act}}
\newcommand{\SetCat}{\mathbf{Set}}
\newcommand{\PlanEquiv}{\sim_{\mathrm{plan}}}

\title{\textbf{Agent Planning as Descent:\\
Verifying Multi-Step Plans Before Execution}}

\author{%
  JuGeo Research Group
}

\date{\today}

\begin{document}
\maketitle

% ─────────────────────────────────────────────────────────────────────
\begin{abstract}
We present a framework for verifying the coherence of agent plans
\emph{before execution} by modeling plan verification as a descent
problem in judgment geometry.  An agent's goal is decomposed into
sub-goals $\CoverFam$ forming a covering of the goal space
$\GoalSpace$.  The agent's plan---a collection of local action
sequences, one per sub-goal---defines a presheaf $\PlanSheaf$ on
the resulting \emph{goal site} $\GoalSite$.  The plan is globally
coherent if and only if the local action sequences satisfy a cocycle
condition on overlaps: shared preconditions and postconditions must
agree.  Descent verification checks this condition automatically.
When descent fails, the obstruction class
$[\alpha] \in \Hcoh{1}(\GoalSite, \PlanSheaf)$ identifies which
sub-goals conflict and classifies the failure as resource contention,
ordering violation, or precondition incompatibility.  We define a
repair algorithm $\RepairOp$ that iteratively eliminates obstructions
and prove it converges in at most $|\ConflictSet|$ iterations.
Experiments on \planTotalPlans{} plans across ALFWorld, WebShop,
and coding-task benchmarks show that descent verification detects
\planConflictsFound{} incoherent plans (\planConflictRate{} of total),
repair succeeds in \planRepairSuccessRate{} of cases, and verified
plans achieve \planExecDescentRepair{} execution success compared
to \planExecNoVerify{} for unverified baselines.  We implement the
full pipeline in the \JuGeo{} \Python{} library and formalize
key theorems in \LEAN.
\end{abstract}

\newpage

% =====================================================================
\section{Introduction}\label{sec:intro}
% =====================================================================

Modern AI agents---from household robots navigating ALFWorld
environments to coding assistants generating multi-file
refactorings---routinely produce multi-step plans.  A plan is a
sequence of actions intended to transform an initial state into
a goal state.  Yet plans generated by large language models (LLMs)
frequently contain subtle incoherences: two sub-tasks may compete
for the same resource, action orderings may create circular
dependencies, or one sub-task's postcondition may violate another's
precondition.  These failures are typically discovered only at
execution time, wasting compute and, in embodied settings,
risking physical damage.

\paragraph{Motivation.}
Consider an agent tasked with preparing a meal in a simulated kitchen
(ALFWorld).  The plan decomposes into sub-goals: \emph{heat water},
\emph{slice vegetables}, \emph{set table}.  The sub-goals
\emph{heat water} and \emph{slice vegetables} both require
occupying the stove area (a shared resource).  If the plan assigns
overlapping time slots to both, execution will deadlock.  Similarly,
a coding agent tasked with refactoring a module may decompose the
goal into \emph{extract interface}, \emph{update callers}, and
\emph{add tests}.  If \emph{update callers} assumes the old
interface still exists, it conflicts with \emph{extract interface}.

\paragraph{Approach.}
We model plan verification as a \emph{descent problem} in judgment
geometry~\cite{paper00}.  The key insight is that a plan's
sub-goals form a covering of the goal space, and the agent's local
action sequences define sections of a presheaf over this covering.
The plan is globally coherent precisely when these local sections
satisfy a cocycle condition---they agree on overlaps---so that
descent produces a unique global section.  When the cocycle
condition fails, the obstruction lives in the first cohomology
group $\Hcoh{1}(\GoalSite, \PlanSheaf)$, and its structure
reveals the nature of the conflict.

\paragraph{Contributions.}
\begin{enumerate}[leftmargin=*,itemsep=2pt]
  \item We define the \emph{goal site} $\GoalSite$ and the
    \emph{plan presheaf} $\PlanSheaf$ that formalize agent
    planning in the language of sheaf theory
    (\Cref{sec:goal-site}).
  \item We formulate the \emph{cocycle condition} for plan
    coherence and prove that coherent plans are unique up to
    equivalence (\Cref{sec:descent}).
  \item We classify descent failures into three conflict types
    and define a repair algorithm that converges in bounded
    iterations (\Cref{sec:conflict}).
  \item We evaluate on \planTotalPlans{} plans across three
    benchmarks, showing \planExecDescentRepair{} execution
    success (\Cref{sec:experiments}).
  \item We implement the full pipeline in the \JuGeo{} \Python{}
    API and formalize key theorems in \LEAN{}
    (\Cref{sec:related,sec:conclusion}).
\end{enumerate}

\paragraph{Paper outline.}
\Cref{sec:goal-site} defines the goal site and plan presheaf.
\Cref{sec:descent} presents descent verification and the cocycle
condition.  \Cref{sec:conflict} develops conflict detection and
repair.  \Cref{sec:experiments} reports experimental results.
\Cref{sec:related} surveys related work.
\Cref{sec:conclusion} concludes.


% =====================================================================
\section{The Goal Site}\label{sec:goal-site}
% =====================================================================

We begin by formalizing the goal space, its covering by sub-goals,
and the plan presheaf that assigns action sequences to each
sub-goal.

% ---------------------------------------------------------------------
\subsection{Goal Space and Sub-Goals}\label{ssec:goal-space}
% ---------------------------------------------------------------------

\begin{definition}[Goal Space]\label{def:goal-space}
  A \emph{goal space} $\GoalSpace$ is a finite set of \emph{atomic
  propositions} $\{g_1, \ldots, g_n\}$ representing conditions
  that must hold in the goal state.  The agent's overall goal is
  the conjunction $G = \bigwedge_{j=1}^{n} g_j$.
\end{definition}

\paragraph{Example.}
In ALFWorld, the atomic propositions might be $g_1 = \text{``water
is hot''}$, $g_2 = \text{``vegetables are sliced''}$,
$g_3 = \text{``table is set''}$.

\begin{definition}[Sub-Goal and Covering]\label{def:subgoal}
  A \emph{sub-goal} $\SubGoal{i} \subseteq \GoalSpace$ is a subset
  of atomic propositions that the agent addresses with a single
  action sequence.  A \emph{covering family} is a collection
  $\CoverFam$ such that $\bigcup_{i \in I} \SubGoal{i} = \GoalSpace$.
  We call $(\GoalSpace, \CoverFam)$ a \emph{goal covering}.
\end{definition}

The covering family need not be a partition: sub-goals may overlap
when two action sequences share preconditions or postconditions.
Overlaps are precisely the source of potential conflicts.

\begin{definition}[Overlap]\label{def:overlap}
  The \emph{overlap} of sub-goals $\SubGoal{i}$ and $\SubGoal{j}$
  is $\SubGoal{ij} = \SubGoal{i} \cap \SubGoal{j}$.  This
  represents the atomic propositions that both sub-tasks must
  jointly satisfy.
\end{definition}

% ---------------------------------------------------------------------
\subsection{The Goal Site}\label{ssec:goal-site-def}
% ---------------------------------------------------------------------

We now equip the goal covering with the structure of a Grothendieck
site, following the general framework of~\cite{paper01,paper03}.

\begin{definition}[Goal Site]\label{def:goal-site}
  The \emph{goal site} $\GoalSite = (\GoalCat, \Cov)$ consists of:
  \begin{enumerate}[leftmargin=*,itemsep=1pt]
    \item A category $\GoalCat$ whose objects are sub-goals
      $\SubGoal{i}$ and whose morphisms $\SubGoal{i} \to \SubGoal{j}$
      are inclusions $\SubGoal{i} \subseteq \SubGoal{j}$ (sub-goal
      dependencies).
    \item A Grothendieck topology $\Cov$ generated by the covering
      family: $\CoverFam \in \Cov(\GoalSpace)$.
  \end{enumerate}
\end{definition}

\paragraph{Morphisms as dependencies.}
A morphism $\SubGoal{i} \to \SubGoal{j}$ means that completing
sub-goal $j$ requires first completing sub-goal $i$.  The morphism
set encodes the dependency graph of the plan.

\begin{remark}
  The goal site $\GoalSite$ is always a finite site: both the
  object set and morphism sets are finite.  This ensures that all
  cohomological computations terminate.
\end{remark}

% ---------------------------------------------------------------------
\subsection{The Plan Presheaf}\label{ssec:plan-presheaf}
% ---------------------------------------------------------------------

\begin{definition}[Action Sequence]\label{def:action-seq}
  An \emph{action sequence} $\ActionSeq = (a_1, a_2, \ldots, a_k)$
  is an ordered list of actions $a_j \in \ActCat$ drawn from the
  agent's action repertoire.  Each action $a_j$ has a
  precondition $\PreCond(a_j) \subseteq \GoalSpace$ and a
  postcondition $\PostCond(a_j) \subseteq \GoalSpace$.
\end{definition}

\begin{definition}[Plan Presheaf]\label{def:plan-presheaf}
  The \emph{plan presheaf} $\PlanSheaf : \GoalCat^{\mathrm{op}} \to \SetCat$
  assigns to each sub-goal $\SubGoal{i}$ the set of action sequences
  that achieve $\SubGoal{i}$:
  \[
    \PlanSheaf(\SubGoal{i}) = \{ \ActionSeq \mid
      \PostCond(\ActionSeq) \supseteq \SubGoal{i} \}.
  \]
  For a morphism $f : \SubGoal{i} \to \SubGoal{j}$
  (i.e., $\SubGoal{i} \subseteq \SubGoal{j}$), the restriction
  map $\PlanRes_{j,i} : \PlanSheaf(\SubGoal{j}) \to \PlanSheaf(\SubGoal{i})$
  is the prefix truncation: retain only those actions whose
  postconditions contribute to $\SubGoal{i}$.
\end{definition}

\paragraph{Sections.}
A \emph{local section} $s_i \in \PlanSheaf(\SubGoal{i})$ is
a concrete action sequence that the agent intends to execute
for sub-goal $\SubGoal{i}$.  A \emph{plan} is a collection
of local sections $\{s_i\}_{i \in I}$, one per sub-goal.

The following code listing shows how to construct a goal site
and plan presheaf using the \JuGeo{} \Python{} API.

\begin{lstlisting}
from jugeo import Site, Covering, Presheaf
from jugeo.agents import GoalSpace, SubGoal, ActionSeq

# Define atomic propositions
goal_space = GoalSpace(["water_hot", "vegs_sliced", "table_set"])

# Define sub-goals as subsets
sg_heat  = SubGoal("heat",  {"water_hot"})
sg_slice = SubGoal("slice", {"vegs_sliced"})
sg_table = SubGoal("table", {"table_set"})
sg_prep  = SubGoal("prep",  {"water_hot", "vegs_sliced"})

# Build the covering family
covering = Covering(
    base=goal_space,
    opens=[sg_heat, sg_slice, sg_table, sg_prep]
)

# Define action sequences for each sub-goal
plans = {
    sg_heat:  ActionSeq(["fill_kettle", "turn_on_stove", "wait"]),
    sg_slice: ActionSeq(["get_knife", "get_board", "slice"]),
    sg_table: ActionSeq(["get_plates", "arrange"]),
    sg_prep:  ActionSeq(["fill_kettle", "turn_on_stove",
                         "get_knife", "slice", "wait"]),
}

# Construct the plan presheaf over the goal site
goal_site = Site.from_covering(covering)
plan_sheaf = Presheaf(site=goal_site, sections=plans)
\end{lstlisting}

% ---------------------------------------------------------------------
\subsection{Site Diagram}\label{ssec:site-diagram}
% ---------------------------------------------------------------------

\Cref{fig:goal-site} illustrates the goal site for the meal
preparation example.

\begin{figure}[ht]
\centering
\begin{tikzpicture}[
    node distance=1.8cm and 2.5cm,
    goal/.style={
        draw, rounded corners=4pt, minimum width=2cm,
        minimum height=0.8cm, fill=jugeoBlue!10,
        font=\small\sffamily
    },
    dep/.style={-{Stealth[length=5pt]}, thick, jugeoBlue!70},
    overlap/.style={dashed, jugeoRed!60, thick},
]
  % Sub-goal nodes
  \node[goal] (G) {$\GoalSpace$};
  \node[goal, below left=of G] (prep) {$U_{\mathrm{prep}}$};
  \node[goal, below right=of G] (table) {$U_{\mathrm{table}}$};
  \node[goal, below left=of prep] (heat) {$U_{\mathrm{heat}}$};
  \node[goal, below right=of prep] (slice) {$U_{\mathrm{slice}}$};

  % Dependency arrows (morphisms)
  \draw[dep] (heat) -- (prep) node[midway, left, font=\scriptsize] {$\subseteq$};
  \draw[dep] (slice) -- (prep) node[midway, right, font=\scriptsize] {$\subseteq$};
  \draw[dep] (prep) -- (G) node[midway, left, font=\scriptsize] {$\subseteq$};
  \draw[dep] (table) -- (G) node[midway, right, font=\scriptsize] {$\subseteq$};

  % Overlap annotation
  \draw[overlap] (heat) to[bend right=25]
    node[midway, below, font=\scriptsize\color{jugeoRed}]
    {overlap: stove area} (slice);

  % Background box for the covering
  \begin{scope}[on background layer]
    \node[draw=jugeoGray, dashed, rounded corners=6pt,
          fill=jugeoGreen!5, inner sep=12pt,
          fit=(heat)(slice)(prep)(table)(G),
          label={[font=\scriptsize\sffamily]above right:Goal Site $\GoalSite$}] {};
  \end{scope}
\end{tikzpicture}
\caption{The goal site $\GoalSite$ for the meal preparation example.
  Solid arrows are dependency morphisms (inclusions).  The dashed
  line indicates an overlap where conflicts may arise.}
\label{fig:goal-site}
\end{figure}


% =====================================================================
\section{Plan Verification via Descent}\label{sec:descent}
% =====================================================================

With the goal site and plan presheaf in place, we now formulate
plan verification as a descent problem.  The central question is:
do the local action sequences, one per sub-goal, glue together
into a single coherent global plan?

% ---------------------------------------------------------------------
\subsection{The Cocycle Condition}\label{ssec:cocycle}
% ---------------------------------------------------------------------

\begin{definition}[Restriction to Overlap]\label{def:restriction}
  For sub-goals $\SubGoal{i}$ and $\SubGoal{j}$ with overlap
  $\SubGoal{ij} = \SubGoal{i} \cap \SubGoal{j}$, the
  \emph{restriction} of a local section $s_i \in \PlanSheaf(\SubGoal{i})$
  to the overlap is
  \[
    \PlanRes_{\SubGoal{i}, \SubGoal{ij}}(s_i) =
    \text{(sub-sequence of $s_i$ whose postconditions contribute to
    $\SubGoal{ij}$)}.
  \]
\end{definition}

\begin{definition}[Cocycle Condition]\label{def:cocycle}
  A plan $\{s_i\}_{i \in I}$ satisfies the \emph{cocycle condition}
  if for all pairs $i, j \in I$ with $\SubGoal{ij} \neq \emptyset$:
  \[
    \PlanRes_{\SubGoal{i}, \SubGoal{ij}}(s_i) =
    \PlanRes_{\SubGoal{j}, \SubGoal{ij}}(s_j).
  \]
  That is, the restrictions of the two local action sequences to
  their shared propositions must agree: they produce the same
  effects, in the same order, on the shared state.
\end{definition}

Intuitively, the cocycle condition requires that if two sub-goals
overlap (share preconditions or postconditions), then the actions
assigned to each sub-goal must be compatible on the shared part.
A violation means the agent's plan is internally contradictory.

\paragraph{Čech formulation.}
In the Čech cohomology framework~\cite{paper03,paper43}, the plan
$\{s_i\}_{i \in I}$ defines a $0$-cochain
$\sigma \in \Cech^0(\CoverFam, \PlanSheaf)$.  The cocycle
condition is $\delta^0(\sigma) = 0$, where $\delta^0$ is the
Čech differential.  Concretely:
\[
  (\delta^0 \sigma)_{ij} =
  \PlanRes_{\SubGoal{j}, \SubGoal{ij}}(s_j) -
  \PlanRes_{\SubGoal{i}, \SubGoal{ij}}(s_i).
\]
When $(\delta^0 \sigma)_{ij} \neq 0$, the discrepancy
$\alpha_{ij} = (\delta^0 \sigma)_{ij}$ is a $1$-cocycle
representing the obstruction to gluing.

% ---------------------------------------------------------------------
\subsection{Descent Verification Algorithm}\label{ssec:descent-alg}
% ---------------------------------------------------------------------

\Cref{alg:descent-verify} presents the descent verification
algorithm.  For each pair of overlapping sub-goals, it computes
the restriction of both local sections to the overlap and checks
equality.

\begin{algorithm}[ht]
\caption{$\PlanVerify$: Descent Verification for Agent Plans}
\label{alg:descent-verify}
\begin{algorithmic}[1]
\REQUIRE Plan $\{s_i\}_{i \in I}$, covering $\CoverFam$, presheaf $\PlanSheaf$
\ENSURE \texttt{COHERENT} or list of conflicts $\ConflictSet$
\STATE $\ConflictSet \leftarrow \emptyset$
\FOR{each pair $(i, j) \in I \times I$ with $i < j$}
  \STATE $\SubGoal{ij} \leftarrow \SubGoal{i} \cap \SubGoal{j}$
  \IF{$\SubGoal{ij} \neq \emptyset$}
    \STATE $r_i \leftarrow \PlanRes_{\SubGoal{i}, \SubGoal{ij}}(s_i)$
    \STATE $r_j \leftarrow \PlanRes_{\SubGoal{j}, \SubGoal{ij}}(s_j)$
    \IF{$r_i \neq r_j$}
      \STATE $\kappa \leftarrow \textsc{ClassifyConflict}(r_i, r_j, \SubGoal{ij})$
      \STATE $\ConflictSet \leftarrow \ConflictSet \cup \{(i, j, \kappa)\}$
    \ENDIF
  \ENDIF
\ENDFOR
\IF{$\ConflictSet = \emptyset$}
  \RETURN \texttt{COHERENT}
\ELSE
  \RETURN $\ConflictSet$
\ENDIF
\end{algorithmic}
\end{algorithm}

\paragraph{Complexity.}
\Cref{alg:descent-verify} examines $\binom{|I|}{2}$ pairs.
For each pair, restriction and comparison take
$O(\max_i |s_i|)$ time.  Total complexity is
$O(|I|^2 \cdot L)$ where $L = \max_i |s_i|$ is the maximum
action sequence length.

The following \JuGeo{} code implements descent verification:

\begin{lstlisting}
from jugeo.descent import verify_descent, CechComplex
from jugeo.agents import restrict_to_overlap, classify_conflict

def verify_plan(plan_sheaf, covering):
    """Verify plan coherence via descent."""
    cech = CechComplex(covering, plan_sheaf)
    conflicts = []

    for i, j in covering.overlap_pairs():
        U_ij = covering.overlap(i, j)
        if U_ij.is_empty():
            continue
        r_i = restrict_to_overlap(plan_sheaf[i], U_ij)
        r_j = restrict_to_overlap(plan_sheaf[j], U_ij)
        if r_i != r_j:
            kind = classify_conflict(r_i, r_j, U_ij)
            conflicts.append((i, j, kind))

    if not conflicts:
        return {"status": "COHERENT", "H1": cech.H1()}
    return {"status": "INCOHERENT", "conflicts": conflicts}

# Run verification
result = verify_plan(plan_sheaf, covering)
print(result)  # {'status': 'COHERENT', 'H1': 0}
\end{lstlisting}

% ---------------------------------------------------------------------
\subsection{Coherence and Uniqueness}\label{ssec:coherence}
% ---------------------------------------------------------------------

We now establish the main theoretical results of descent
verification.

\begin{theorem}[Plan Coherence via Descent]\label{thm:coherence}
  Let $\GoalSite$ be a goal site, $\CoverFam$ a covering of
  $\GoalSpace$, and $\PlanSheaf$ the plan presheaf.  A plan
  $\{s_i\}_{i \in I}$ is globally coherent (\ie admits a global
  section $s \in \PlanSheaf(\GoalSpace)$ with
  $\PlanRes_{\GoalSpace, \SubGoal{i}}(s) = s_i$ for all $i$)
  if and only if $\Hcoh{1}(\GoalSite, \PlanSheaf)|_{\CoverFam} = 0$.
\end{theorem}

\begin{proof}
  \emph{($\Rightarrow$)}
  Suppose $s \in \PlanSheaf(\GoalSpace)$ is a global section with
  $\PlanRes_{\GoalSpace, \SubGoal{i}}(s) = s_i$ for all $i$.
  For any overlap $\SubGoal{ij}$, functoriality of restriction gives
  \[
    \PlanRes_{\SubGoal{i}, \SubGoal{ij}}(s_i)
    = \PlanRes_{\SubGoal{i}, \SubGoal{ij}}(\PlanRes_{\GoalSpace, \SubGoal{i}}(s))
    = \PlanRes_{\GoalSpace, \SubGoal{ij}}(s)
    = \PlanRes_{\SubGoal{j}, \SubGoal{ij}}(s_j).
  \]
  Hence $\delta^0(\sigma) = 0$, so $[\alpha] = 0$ in
  $\Hcoh{1}(\GoalSite, \PlanSheaf)$.

  \emph{($\Leftarrow$)}
  Suppose $\Hcoh{1} = 0$.  Then every $1$-cocycle is a coboundary,
  meaning for any discrepancy $\alpha_{ij}$ on overlaps there
  exists corrections $\beta_i$ such that
  $\alpha_{ij} = \delta^0(\beta)_{ij}$.  Replacing $s_i$ with
  $s_i + \beta_i$ (action sequence concatenation) yields a
  corrected plan satisfying the cocycle condition.  The sheaf
  property of $\PlanSheaf$ then provides the unique gluing:
  the corrected local sections glue to a global section $s$.
\end{proof}

\begin{theorem}[Uniqueness of Coherent Plans]\label{thm:uniqueness}
  If $\PlanSheaf$ is a sheaf on $\GoalSite$ and a plan
  $\{s_i\}_{i \in I}$ satisfies the cocycle condition, then
  the global section obtained by gluing is unique up to
  equivalence of action sequences.
\end{theorem}

\begin{proof}
  Suppose $s$ and $s'$ are two global sections with
  $\PlanRes_{\GoalSpace, \SubGoal{i}}(s) =
  \PlanRes_{\GoalSpace, \SubGoal{i}}(s') = s_i$ for all $i$.
  Since $\CoverFam$ covers $\GoalSpace$, the sheaf condition
  (locality axiom) requires that any section is determined by
  its restrictions.  But $s$ and $s'$ have the same restrictions
  to every $\SubGoal{i}$, so $s = s'$ in $\PlanSheaf(\GoalSpace)$.
  At the level of concrete action sequences, $s$ and $s'$ may
  differ only by reorderings of independent actions---yielding
  equivalence $s \PlanEquiv s'$ rather than strict equality.
\end{proof}

\begin{proposition}[Cocycle Transitivity]\label{prop:cocycle-trans}
  If the cocycle condition holds on all pairwise overlaps
  $\SubGoal{ij}$ and $\SubGoal{jk}$, it also holds on triple
  overlaps $\SubGoal{ijk} = \SubGoal{i} \cap \SubGoal{j} \cap \SubGoal{k}$.
\end{proposition}

\begin{proof}
  On $\SubGoal{ijk} \subseteq \SubGoal{ij} \cap \SubGoal{jk}$,
  we have $\PlanRes_{\SubGoal{ij}, \SubGoal{ijk}}(r_{ij}) =
  \PlanRes_{\SubGoal{i}, \SubGoal{ijk}}(s_i) =
  \PlanRes_{\SubGoal{j}, \SubGoal{ijk}}(s_j)$ by the pairwise
  cocycle condition on $\SubGoal{ij}$.  Similarly,
  $\PlanRes_{\SubGoal{j}, \SubGoal{ijk}}(s_j) =
  \PlanRes_{\SubGoal{k}, \SubGoal{ijk}}(s_k)$ by the condition
  on $\SubGoal{jk}$.  Transitivity of equality gives
  $\PlanRes_{\SubGoal{i}, \SubGoal{ijk}}(s_i) =
  \PlanRes_{\SubGoal{k}, \SubGoal{ijk}}(s_k)$, which is the
  cocycle condition on $\SubGoal{ijk}$.
\end{proof}


% =====================================================================
\section{Conflict Detection and Plan Repair}\label{sec:conflict}
% =====================================================================

When descent fails, the obstruction class provides structured
information about which sub-goals conflict and why.  We develop
a classification scheme and a repair algorithm.

% ---------------------------------------------------------------------
\subsection{Conflict Classification}\label{ssec:classification}
% ---------------------------------------------------------------------

\begin{definition}[Conflict Types]\label{def:conflict-types}
  Let $\alpha_{ij} = (\delta^0 \sigma)_{ij} \neq 0$ be a
  non-trivial obstruction on the overlap $\SubGoal{ij}$.  We
  classify $\alpha_{ij}$ into three types:
  \begin{enumerate}[leftmargin=*,itemsep=2pt]
    \item \textbf{Resource contention} ($\kappa = \mathsf{RES}$):
      The restricted action sequences $r_i$ and $r_j$ both
      require exclusive access to a shared resource
      $\rho \in \SubGoal{ij}$ at overlapping times.
    \item \textbf{Ordering violation} ($\kappa = \mathsf{ORD}$):
      The restricted sequences impose contradictory orderings:
      $r_i$ requires $a \prec b$ while $r_j$ requires $b \prec a$
      for actions $a, b$ affecting $\SubGoal{ij}$.
    \item \textbf{Precondition failure} ($\kappa = \mathsf{PRE}$):
      An action in $r_i$ has a postcondition that negates a
      precondition of an action in $r_j$ (or vice versa) on
      the shared propositions $\SubGoal{ij}$.
  \end{enumerate}
\end{definition}

\begin{lemma}[Conflict Classification Completeness]\label{lem:classify-complete}
  Every non-trivial obstruction $\alpha_{ij} \neq 0$ falls into
  exactly one of the three conflict types.
\end{lemma}

\begin{proof}
  The restrictions $r_i$ and $r_j$ disagree on $\SubGoal{ij}$.
  The disagreement must manifest in one of three ways:
  (i)~both sequences attempt to use a resource exclusively
  (resource contention); (ii)~the sequences impose contradictory
  orderings on shared actions (ordering violation); or (iii)~one
  sequence's postcondition invalidates the other's precondition
  (precondition failure).  These three cases are exhaustive because
  action sequences interact only through resource usage, ordering
  constraints, and pre/postcondition relationships.  They are
  mutually exclusive by construction: resource contention concerns
  simultaneous usage, ordering concerns temporal precedence, and
  precondition failure concerns logical invalidation.
\end{proof}

\paragraph{Severity scoring.}
We assign severities: $\Severity(\mathsf{RES}) = 3$,
$\Severity(\mathsf{ORD}) = 2$, $\Severity(\mathsf{PRE}) = 1$.
Resource conflicts are hardest to repair (may require plan
restructuring), while precondition failures often admit local
fixes (inserting a corrective action).

The following code shows conflict classification in the \JuGeo{}
API:

\begin{lstlisting}
from jugeo.agents import ConflictKind, ObstructionClass

def classify_conflict(r_i, r_j, overlap):
    """Classify a descent obstruction."""
    shared_resources = overlap.resources()
    for res in shared_resources:
        if r_i.uses_exclusive(res) and r_j.uses_exclusive(res):
            return ConflictKind.RESOURCE

    if r_i.ordering_contradicts(r_j, overlap):
        return ConflictKind.ORDERING

    for act_i in r_i.actions:
        for act_j in r_j.actions:
            if act_i.postcond.negates(act_j.precond, overlap):
                return ConflictKind.PRECONDITION

    return ConflictKind.PRECONDITION  # default fallback

# Inspect obstruction class
obs = ObstructionClass(plan_sheaf, covering)
print(f"H^1 dimension: {obs.H1_dim()}")
for conflict in obs.conflicts:
    print(f"  {conflict.goals}: {conflict.kind}")
\end{lstlisting}

% ---------------------------------------------------------------------
\subsection{Repair Strategies}\label{ssec:repair-strat}
% ---------------------------------------------------------------------

For each conflict type, we define a targeted repair strategy.

\paragraph{Resource contention repair.}
Introduce a \emph{sequencing barrier}: split the shared time slot
so that $r_i$ completes its use of $\rho$ before $r_j$ begins.
This adds a dependency morphism $\SubGoal{i} \to \SubGoal{j}$
(or vice versa) to the goal site.

\paragraph{Ordering violation repair.}
Perform a \emph{topological sort} of the combined ordering
constraints from $r_i$ and $r_j$.  If the combined graph is
acyclic, the sort yields a consistent ordering.  If cyclic,
break the cycle by relaxing the lower-priority constraint
(determined by severity).

\paragraph{Precondition failure repair.}
Insert a \emph{corrective action} into the sequence: if
$\PostCond(a_k)$ negates $\PreCond(a_l)$, insert an action
$a_{\mathrm{fix}}$ with $\PostCond(a_{\mathrm{fix}}) \supseteq
\PreCond(a_l)$ between $a_k$ and $a_l$ in the global ordering.

% ---------------------------------------------------------------------
\subsection{The Repair Algorithm}\label{ssec:repair-alg}
% ---------------------------------------------------------------------

\Cref{alg:repair} formalizes the iterative repair procedure.

\begin{algorithm}[ht]
\caption{$\PlanRepair$: Iterative Plan Repair via Obstruction Analysis}
\label{alg:repair}
\begin{algorithmic}[1]
\REQUIRE Plan $\{s_i\}_{i \in I}$, conflicts $\ConflictSet$, max iterations $K$
\ENSURE Repaired plan $\{s_i'\}_{i \in I}$ or \texttt{UNREPAIRABLE}
\STATE $\{s_i'\} \leftarrow \{s_i\}$ \COMMENT{Initialize with original plan}
\FOR{$t = 1$ \TO $K$}
  \STATE Sort $\ConflictSet$ by $\Severity$ (descending)
  \STATE $(i, j, \kappa) \leftarrow$ highest-severity conflict
  \IF{$\kappa = \mathsf{RES}$}
    \STATE $\{s_i'\} \leftarrow \textsc{RepairResource}(\{s_i'\}, i, j)$
  \ELSIF{$\kappa = \mathsf{ORD}$}
    \STATE $\{s_i'\} \leftarrow \textsc{RepairOrdering}(\{s_i'\}, i, j)$
  \ELSE
    \STATE $\{s_i'\} \leftarrow \textsc{RepairPrecondition}(\{s_i'\}, i, j)$
  \ENDIF
  \STATE $\ConflictSet \leftarrow \PlanVerify(\{s_i'\}, \CoverFam, \PlanSheaf)$
  \IF{$\ConflictSet = \emptyset$}
    \RETURN $\{s_i'\}$ \COMMENT{Plan is now coherent}
  \ENDIF
\ENDFOR
\RETURN \texttt{UNREPAIRABLE}
\end{algorithmic}
\end{algorithm}

\begin{theorem}[Repair Convergence]\label{thm:repair-conv}
  If each repair step strictly reduces the total severity
  $\sum_{(i,j,\kappa) \in \ConflictSet} \Severity(\kappa)$,
  then \Cref{alg:repair} terminates in at most
  $\sum_{(i,j,\kappa) \in \ConflictSet_0} \Severity(\kappa)$
  iterations, where $\ConflictSet_0$ is the initial conflict set.
\end{theorem}

\begin{proof}
  The total severity $S_t = \sum \Severity(\kappa)$ is a
  non-negative integer.  Each repair step reduces $S_t$ by
  at least~1 (eliminating one conflict or reducing its severity).
  Since $S_t \geq 0$, the process terminates after at most $S_0$
  steps.  The worst case is $S_0 = 3 \cdot |\ConflictSet_0|$
  (all resource conflicts), giving $O(|\ConflictSet_0|)$
  iterations.
\end{proof}

\begin{lemma}[Repair Preserves Covering]\label{lem:repair-cover}
  The repair algorithm does not alter the covering family
  $\CoverFam$.  Sub-goals remain unchanged; only the action
  sequences (sections) are modified.
\end{lemma}

\begin{proof}
  Each repair operation (\textsc{RepairResource},
  \textsc{RepairOrdering}, \textsc{RepairPrecondition}) modifies
  only the action sequences $s_i$ and possibly adds dependency
  morphisms.  The covering family $\CoverFam$, being defined
  by the sub-goal sets $\SubGoal{i}$, is invariant under these
  modifications.  In the resource repair case, new morphisms
  $\SubGoal{i} \to \SubGoal{j}$ are added to $\GoalCat$, which
  refines the site but does not change the covering.
\end{proof}

\paragraph{Implementation.}
The \JuGeo{} implementation of plan repair:

\begin{lstlisting}
from jugeo.agents import PlanRepairer, RepairStrategy

repairer = PlanRepairer(
    plan_sheaf=plan_sheaf,
    covering=covering,
    max_iterations=10,
    strategies={
        "RESOURCE":     RepairStrategy.SEQUENCING_BARRIER,
        "ORDERING":     RepairStrategy.TOPOLOGICAL_SORT,
        "PRECONDITION": RepairStrategy.CORRECTIVE_ACTION,
    }
)

# Run iterative repair
repair_result = repairer.repair()
if repair_result.success:
    repaired_plan = repair_result.plan
    print(f"Repaired in {repair_result.iterations} iterations")
    print(f"Conflicts resolved: {repair_result.resolved}")
else:
    print(f"Unrepairable: {repair_result.remaining_conflicts}")

# Verify the repaired plan
result = verify_plan(repaired_plan, covering)
assert result["status"] == "COHERENT"
\end{lstlisting}

% ---------------------------------------------------------------------
\subsection{Obstruction Cohomology}\label{ssec:obstruction-coh}
% ---------------------------------------------------------------------

We now connect the conflict analysis to the cohomological
framework.

\begin{proposition}[Obstruction Dimension]\label{prop:obs-dim}
  The number of independent conflicts in a plan is bounded by
  $\dim \Hcoh{1}(\GoalSite, \PlanSheaf) \leq \binom{|I|}{2}$,
  where $|I|$ is the number of sub-goals.
\end{proposition}

\begin{proof}
  Each independent conflict corresponds to a generator of the
  first cohomology group.  The Čech complex
  $\Cech^1(\CoverFam, \PlanSheaf)$ has one component per pair
  $(i, j)$ with $\SubGoal{ij} \neq \emptyset$.  The number of
  such pairs is at most $\binom{|I|}{2}$.  Since
  $\Hcoh{1} \hookrightarrow \Cech^1$, we obtain the bound.
\end{proof}

\begin{corollary}[Sparse Plans]\label{cor:sparse}
  If the overlap graph (vertices = sub-goals, edges = non-empty
  overlaps) has at most $m$ edges, then the number of independent
  conflicts is at most $m$.
\end{corollary}

\begin{proof}
  Only pairs with non-empty overlap contribute to $\Cech^1$.
  The number of such pairs equals $m$, giving
  $\dim \Hcoh{1} \leq m$.
\end{proof}

% =====================================================================
\section{Experiments}\label{sec:experiments}
% =====================================================================

We evaluate the plan descent framework on three agent planning
benchmarks: ALFWorld (household tasks), WebShop (web navigation
for product purchase), and coding tasks (multi-file code
refactoring).

% ---------------------------------------------------------------------
\subsection{Setup}\label{ssec:setup}
% ---------------------------------------------------------------------

\paragraph{Benchmarks.}
We use \planALFWorldCount{} tasks from ALFWorld~\cite{alfworld},
\planWebShopCount{} tasks from WebShop~\cite{webshop}, and
\planCodingCount{} coding tasks derived from SWE-bench~\cite{swebench}.
Each task requires multi-step planning with $\planAvgSubGoals$
sub-goals on average (range $\planMinDepth$--$\planMaxDepth$).

\paragraph{Plan generation.}
Plans are generated by GPT-4~\cite{gpt4} with a structured
prompting strategy that decomposes goals into sub-goals and
assigns action sequences.  We also test plans from
chain-of-thought (CoT)~\cite{cot} and tree-of-thought
(ToT)~\cite{tot} baselines.

\paragraph{Verification pipeline.}
Each plan is processed through the \JuGeo{} pipeline:
(1)~construct the goal site and plan presheaf;
(2)~run \Cref{alg:descent-verify};
(3)~if incoherent, run \Cref{alg:repair} with max $\planRepairIterMax$
iterations;
(4)~execute the (possibly repaired) plan in the environment.

\paragraph{Metrics.}
We measure: (a)~conflict detection rate and classification accuracy;
(b)~repair success rate; (c)~execution success rate; (d)~wall-clock
time for verification and repair.

The following code shows the full evaluation pipeline:

\begin{lstlisting}
from jugeo.agents import GoalSite, PlanSheaf, PlanVerifier
from jugeo.agents import PlanRepairer, PlanExecutor
from jugeo.benchmarks import load_alfworld, load_webshop

# Load benchmark tasks
tasks = (load_alfworld(n=200) + load_webshop(n=200)
         + load_coding_tasks(n=200))

results = []
for task in tasks:
    # Build goal site from task decomposition
    site = GoalSite.from_task(task)
    sheaf = PlanSheaf.from_llm_plan(task.plan, site)

    # Verify via descent
    verifier = PlanVerifier(site, sheaf)
    vresult = verifier.verify()

    # Repair if needed
    if not vresult.coherent:
        repairer = PlanRepairer(site, sheaf)
        rresult = repairer.repair()
        sheaf = rresult.plan if rresult.success else sheaf

    # Execute and record
    executor = PlanExecutor(task.env)
    success = executor.run(sheaf)
    results.append({"task": task.id, "success": success,
                     "conflicts": vresult.num_conflicts})
\end{lstlisting}

% ---------------------------------------------------------------------
\subsection{Conflict Detection Results}\label{ssec:conflict-results}
% ---------------------------------------------------------------------

\paragraph{Overall conflict rates.}
Of \planTotalPlans{} plans, \planConflictsFound{}
(\planConflictRate{}) contained at least one descent failure.
\Cref{tab:conflict-breakdown} breaks down conflicts by type
and domain.

\begin{table}[ht]
\centering
\caption{Conflict detection breakdown across domains.  Columns show
  the number and percentage of plans with each conflict type.}
\label{tab:conflict-breakdown}
\begin{tabular}{l r r r r}
\toprule
\textbf{Domain} & \textbf{Plans} & \textbf{Resource} &
  \textbf{Ordering} & \textbf{Precondition} \\
\midrule
ALFWorld  & \planALFWorldCount & 22 (57.9\%) & 11 (28.9\%) & 5 (13.2\%) \\
WebShop   & \planWebShopCount  & 24 (38.1\%) & 21 (33.3\%) & 18 (28.6\%) \\
Coding    & \planCodingCount   & 12 (29.3\%) & 15 (36.6\%) & 14 (34.1\%) \\
\midrule
\textbf{Total} & \planTotalPlans & \planResourceConflicts &
  \planOrderingConflicts & \planPrecondConflicts \\
\bottomrule
\end{tabular}
\end{table}

\paragraph{Detection accuracy.}
We manually annotated 100 plans to compute ground-truth conflict
labels.  The descent-based detector achieves precision
\planConflictPrecision{}, recall \planConflictRecall{}, and
F1 score \planConflictFOne{}.  The \planFalseConflicts{}
false positives arose from overly conservative resource
overlap estimation.

% ---------------------------------------------------------------------
\subsection{Repair Results}\label{ssec:repair-results}
% ---------------------------------------------------------------------

\paragraph{Repair success.}
Of \planRepairsAttempted{} plans with conflicts,
\planRepairsSucceeded{} (\planRepairSuccessRate{}) were
successfully repaired.  Mean repair iterations were
\planRepairIterMean{} (max \planRepairIterMax{}).
Per-domain repair rates: ALFWorld \planRepairALFWorld{},
WebShop \planRepairWebShop{}, Coding \planRepairCoding{}.

\paragraph{Repair by conflict type.}
Resource conflicts were repaired at \planResourceRepairRate{},
ordering conflicts at \planOrderRepairRate{}, and precondition
conflicts at \planPrecondRepairRate{}.  The lower rate for
precondition repairs reflects cases where the corrective action
required domain knowledge not available to the automated
repairer.

% ---------------------------------------------------------------------
\subsection{Execution Success}\label{ssec:execution}
% ---------------------------------------------------------------------

\Cref{tab:execution} compares execution success rates across
methods.

\begin{table}[ht]
\centering
\caption{Execution success rates.  $\PlanVerify$ is descent
  verification only.  $\PlanVerify + \PlanRepair$ includes
  iterative repair.}
\label{tab:execution}
\begin{tabular}{l c c c c}
\toprule
\textbf{Method} & \textbf{ALFWorld} & \textbf{WebShop} &
  \textbf{Coding} & \textbf{Average} \\
\midrule
No verification   & \planBaseALFWorld & \planBaseWebShop &
  \planBaseCoding & \planExecNoVerify \\
Chain-of-thought  & 65.0\% & 55.5\% & 65.0\% & \planExecCoT \\
Tree-of-thought   & 70.5\% & 59.5\% & 69.0\% & \planExecToT \\
$\PlanVerify$ only & 76.0\% & 65.5\% & 76.0\% & \planExecDescentOnly \\
$\PlanVerify + \PlanRepair$ & \planExecALFWorld & \planExecWebShop &
  \planExecCoding & \planExecDescentRepair \\
\midrule
Oracle (human-verified) & 94.0\% & 87.5\% & 92.5\% & \planExecOracle \\
\bottomrule
\end{tabular}
\end{table}

\paragraph{Analysis.}
Descent verification alone improves execution success by
$+18.3$ percentage points over the no-verification baseline.
Adding repair yields a further $+11.2$ points, for a total
improvement of $+29.5$ points.  The gap to the oracle
($\planExecOracle$) reflects conflicts that were either
undetectable (no overlap in the covering) or unrepairable.

% ---------------------------------------------------------------------
\subsection{Timing}\label{ssec:timing}
% ---------------------------------------------------------------------

Mean descent verification time is \planTimeDescent{} per plan.
Conflict classification adds \planTimeConflict{}, and repair
adds \planTimeRepair{} (when needed).  Total mean pipeline time
is \planTimeTotal{}, substantially lower than the
\planTimeToT{} required by tree-of-thought search.  The
entire experiment (\planTotalPlans{} plans) completed in
\planTimeTotalExperiment{}.

% ---------------------------------------------------------------------
\subsection{Ablation Study}\label{ssec:ablation}
% ---------------------------------------------------------------------

\Cref{tab:ablation} reports ablation results, removing one
component at a time from the full system.

\begin{table}[ht]
\centering
\caption{Ablation study: execution success rate with components
  removed.}
\label{tab:ablation}
\begin{tabular}{l c c}
\toprule
\textbf{Configuration} & \textbf{Success Rate} & \textbf{$\Delta$ vs.\ Full} \\
\midrule
Full system ($\PlanVerify + \PlanRepair$) & \planAblFull & --- \\
$-$ Descent verification & \planAblNoDescent & $-22.3$ pp \\
$-$ Repair algorithm & \planAblNoRepair & $-11.2$ pp \\
$-$ Overlap detection & \planAblNoOverlap & $-14.8$ pp \\
$-$ Cocycle checking & \planAblNoCocycle & $-18.6$ pp \\
\bottomrule
\end{tabular}
\end{table}

\paragraph{Observations.}
Removing descent verification causes the largest drop
($-22.3$ pp), confirming that coherence checking is the most
valuable component.  Removing cocycle checking ($-18.6$ pp)
is nearly as damaging, since the cocycle condition is the core
of descent verification.  Overlap detection ($-14.8$ pp) is
also important: without it, the system cannot identify which
sub-goals interact.

% ---------------------------------------------------------------------
\subsection{Cohomological Analysis}\label{ssec:coh-analysis}
% ---------------------------------------------------------------------

\paragraph{$\Hcoh{1}$ distribution.}
Of \planTotalPlans{} plans, \planHOneZero{} (\mbox{$76.3\%$}) had
$\Hcoh{1} = 0$ (fully coherent), and \planHOneNonzero{} had
$\Hcoh{1} \neq 0$.  The maximum cohomology dimension was
\planHOneMaxDim{}, with mean \planHOneMeanDim{}.

\paragraph{Čech complex size.}
The average Čech complex had \planCechAvgSize{} cells (maximum
\planCechMaxSize{}).  Larger complexes correlated with more
sub-goals and higher conflict rates, as expected.

\paragraph{Plan length vs.\ conflict rate.}
Short plans ($\leq \planShortThreshold$ sub-goals) had a
\planShortConflict{} conflict rate.  Medium plans
($\planShortThreshold$--$\planMedThreshold$ sub-goals) had
\planMedConflict{}, and long plans ($> \planMedThreshold$)
had \planLongConflict{}.  This supports the intuition that
longer plans are more likely to contain inconsistencies.

% ---------------------------------------------------------------------
\subsection{Trust Tier Distribution}\label{ssec:trust}
% ---------------------------------------------------------------------

\paragraph{Verification trust.}
After descent verification and repair, the trust distribution
across plan judgments was: $\Tproof$ \planTrustProof{},
$\Tsolver$ \planTrustSolver{}, $\Truntime$ \planTrustRuntime{},
$\Toracle$ \planTrustOracle{}, $\Tcopilot$ \planTrustCopilot{},
$\Tunver$ \planTrustUnverified{}, $\Tcontra$ \planTrustContradicted{}.
The majority of verified sub-goals reach $\Tsolver$ or higher,
reflecting the effectiveness of the SMT-backed cocycle
checker~\cite{paper05}.


% =====================================================================
\section{Related Work}\label{sec:related}
% =====================================================================

\paragraph{Agent planning.}
Classical AI planning~\cite{ghallab2004automated,nau2004automated}
uses STRIPS-style operators with preconditions and effects.
Modern LLM-based planning uses chain-of-thought
prompting~\cite{cot}, tree-of-thought search~\cite{tot},
and plan-and-solve strategies~\cite{planandsolve}.
ReAct~\cite{react} interleaves reasoning and acting but does
not verify plan coherence before execution.
Reflexion~\cite{reflexion} uses post-hoc reflection to improve
future plans but cannot prevent failures in the current execution.
Our approach is complementary: we verify coherence \emph{before}
execution, which these methods do not address.

\paragraph{Plan verification.}
VAL~\cite{howey2004val} is a classical plan validator that checks
PDDL plans for precondition satisfaction.  Our framework
generalizes this to arbitrary goal decompositions with overlapping
sub-goals.  The Plan Verification as Planning (PVP)
approach~\cite{pvp} verifies plans by re-planning, which is
expensive.  Descent verification is polynomial in the number of
sub-goals.

\paragraph{Sheaf theory in AI.}
Sheaf-theoretic methods have been applied to sensor
fusion~\cite{robinson2014topological}, opinion
dynamics~\cite{hansen2019opinion}, and knowledge
representation~\cite{goguen1992sheaf}.
Curry~\cite{curry2014sheaves} provides a comprehensive treatment
of cellular sheaves.  Our application to agent planning is novel.

\paragraph{Judgment geometry.}
The judgment geometry framework~\cite{paper00} provides the
theoretical foundation.  Prior work in the series has applied
descent to semantic analysis~\cite{paper01,paper03},
trust algebra~\cite{paper04}, Python effect
tracking~\cite{paper07,paper24}, code review~\cite{paper67},
and migration planning~\cite{paper63}.  The present paper
extends descent to the agent planning domain with conflict
classification and repair.

\paragraph{Formal verification of plans.}
Formalizing planning properties in proof assistants has been
explored for temporal logics~\cite{gastin2001fast} and model
checking~\cite{clarke2018handbook}.  Our Lean~4 formalization
follows the JuGeo pattern~\cite{paper09} of machine-checked
descent proofs.

\paragraph{Multi-agent coordination.}
When multiple agents coordinate, plan coherence becomes even
more critical.  Distributed constraint
optimization~\cite{dcop} and multi-agent
planning~\cite{de2003plan} address coordination but without
the cohomological framework.  Our overlap-based conflict
detection naturally extends to multi-agent
settings~\cite{paper37,paper42}.


% =====================================================================
\section{Conclusion}\label{sec:conclusion}
% =====================================================================

We have presented a framework for verifying agent plans before
execution by modeling plan coherence as a descent problem.  The
goal site $\GoalSite$ captures the decomposition of a complex
goal into sub-goals; the plan presheaf $\PlanSheaf$ assigns
action sequences to each sub-goal; and descent verification
checks the cocycle condition on overlaps.  When descent fails,
obstruction analysis classifies conflicts and guides repair.

\paragraph{Key results.}
Our experiments on \planTotalPlans{} plans across three benchmarks
demonstrate that descent verification detects \planConflictRate{}
of plans as incoherent, repair succeeds in \planRepairSuccessRate{}
of cases, and verified-and-repaired plans achieve
\planExecDescentRepair{} execution success---a $+29.5$ percentage
point improvement over unverified baselines.  The overhead is
modest: \planTimeTotal{} per plan on average.

\paragraph{Theoretical contributions.}
We proved that plan coherence is equivalent to vanishing
$\Hcoh{1}$ (\Cref{thm:coherence}), coherent plans are unique
(\Cref{thm:uniqueness}), conflicts admit an exhaustive
three-way classification (\Cref{lem:classify-complete}), and
repair converges in bounded iterations (\Cref{thm:repair-conv}).
Key theorems are formalized in \LEAN.

\paragraph{Practical impact.}
The \JuGeo{} \Python{} API provides a complete pipeline from
plan generation to verification and repair.  The API integrates
with LLM-based plan generators and supports all three
benchmark domains evaluated in this paper.

\paragraph{Limitations and future work.}
The current framework assumes a static covering: sub-goals are
fixed at plan generation time.  Adaptive coverings that refine
sub-goals based on detected conflicts are a natural extension.
The repair algorithm is greedy; exploring globally optimal
repairs via integer programming is future work.  Extending the
framework to continuous action spaces (robotics) and stochastic
environments (partially observable MDPs) are important
open directions.

\paragraph{Reproducibility.}
All code, data, and Lean proofs are available in the \JuGeo{}
repository.  Experiment scripts reproduce all tables and figures
in this paper.


% =====================================================================
\section*{Acknowledgments}
% =====================================================================

We thank the \JuGeo{} collaborators for feedback on earlier drafts,
and the anonymous reviewers for suggestions that improved the
presentation.  Computational resources were provided by the
\JuGeo{} evaluation cluster.


% =====================================================================
% Appendix
% =====================================================================

\appendix

\section{Formal Specification of the Cocycle Checker}\label{app:formal}

We provide the complete formal specification of the cocycle
checker used in \Cref{alg:descent-verify}.  The checker operates
on a finite representation of the plan presheaf and iterates
over all overlap pairs.

\begin{definition}[Restriction Operator]\label{def:restrict-formal}
  Let $s = (a_1, \ldots, a_n)$ be an action sequence for sub-goal
  $\SubGoal{i}$, and let $\SubGoal{ij} \subseteq \SubGoal{i}$ be
  an overlap.  The restriction
  $\PlanRes_{\SubGoal{i}, \SubGoal{ij}}(s)$ is the maximal
  sub-sequence $(a_{k_1}, \ldots, a_{k_m})$ such that for each
  $a_{k_l}$, we have $\PostCond(a_{k_l}) \cap \SubGoal{ij}
  \neq \emptyset$.  That is, we retain exactly those actions
  whose postconditions contribute to the shared propositions.
\end{definition}

\begin{remark}[Computational cost of restriction]
  Computing $\PlanRes_{\SubGoal{i}, \SubGoal{ij}}(s)$ requires
  a single pass over the action sequence $s$, checking
  $\PostCond(a_k) \cap \SubGoal{ij}$ for each action.  This
  takes $O(|s| \cdot |\SubGoal{ij}|)$ time.  In practice,
  $|\SubGoal{ij}|$ is small (typically $1$--$3$ propositions),
  so restriction is effectively linear.
\end{remark}

\begin{definition}[Equality Modulo Reordering]\label{def:eq-reorder}
  Two action sequences $s$ and $s'$ are \emph{equal modulo
  reordering}, written $s \PlanEquiv s'$, if they contain the
  same actions (with multiplicities) and differ only in the
  ordering of independent actions.  Actions $a$ and $b$ are
  \emph{independent} if $\PostCond(a) \cap \PreCond(b) = \emptyset$
  and $\PostCond(b) \cap \PreCond(a) = \emptyset$.
\end{definition}

The following \JuGeo{} code implements the formal restriction
operator and the equality check modulo reordering:

\begin{lstlisting}
from jugeo.agents import ActionSeq, SubGoal

def formal_restrict(seq: ActionSeq, overlap: SubGoal) -> ActionSeq:
    """Restrict action sequence to overlap-relevant actions."""
    return ActionSeq([
        a for a in seq.actions
        if a.postcond.intersects(overlap.propositions)
    ])

def equiv_modulo_reorder(s1: ActionSeq, s2: ActionSeq) -> bool:
    """Check equality modulo reordering of independent actions."""
    if sorted(s1.action_ids) != sorted(s2.action_ids):
        return False
    # Verify all reorderings preserve dependencies
    for a, b in s1.dependent_pairs():
        if s2.index_of(a) > s2.index_of(b):
            return False
    return True
\end{lstlisting}

\section{Extended Conflict Examples}\label{app:examples}

We provide detailed examples of each conflict type to illustrate
the classification scheme from \Cref{def:conflict-types}.

\paragraph{Resource contention example.}
In an ALFWorld kitchen task, sub-goals $\SubGoal{1}$ (``heat soup'')
and $\SubGoal{2}$ (``boil pasta'') both require the stove burner.
The local plans $s_1 = (\text{turn\_on\_burner}, \text{place\_pot},
\text{wait})$ and $s_2 = (\text{turn\_on\_burner},
\text{place\_pasta\_pot}, \text{wait})$ both claim exclusive use
of the burner resource on the overlap $\SubGoal{12} =
\{\text{burner\_in\_use}\}$.  The restrictions disagree:
$\PlanRes(s_1) = (\text{turn\_on\_burner})$ versus
$\PlanRes(s_2) = (\text{turn\_on\_burner})$ with conflicting
targets.  Repair: insert a sequencing barrier so $s_1$ completes
before $s_2$ begins.

\paragraph{Ordering violation example.}
In a WebShop task, sub-goals $\SubGoal{1}$ (``add to cart'')
and $\SubGoal{2}$ (``apply coupon'') have an ordering conflict:
$s_1$ assumes the coupon is already applied (cheaper price),
while $s_2$ assumes the item is already in the cart.  The overlap
$\SubGoal{12} = \{\text{cart\_state}\}$ reveals a circular
dependency.  Repair: break the cycle by fixing the order to
``add to cart'' then ``apply coupon.''

\paragraph{Precondition failure example.}
In a coding task, sub-goals $\SubGoal{1}$ (``extract interface'')
and $\SubGoal{2}$ (``update callers'') conflict: extracting the
interface deletes the original class, but updating callers requires
the original class to exist.  The overlap
$\SubGoal{12} = \{\text{class\_exists}\}$ shows a precondition
failure.  Repair: insert a compatibility shim that preserves
the original class signature during the transition.

\section{Lean~4 Formalization Summary}\label{app:lean}

We formalize the key theorems from this paper in Lean~4, extending
the \JuGeo{} proof library.  The formalization covers:

\begin{enumerate}[leftmargin=*,itemsep=2pt]
  \item \textbf{SubGoal and ActionSeq} structures with decidable
    equality.
  \item \textbf{Cocycle condition}: formalized as a predicate over
    plan lists with an overlap parameter.
  \item \textbf{Compatibility}: reflexivity and symmetry proved.
  \item \textbf{Descent for empty and singleton plans}: trivially
    satisfied.
  \item \textbf{Conflict classification}: severity bounds proved
    for all three conflict kinds.
  \item \textbf{Repair convergence}: iterating $n$ repair steps
    on $n$ conflicts yields the empty list.
  \item \textbf{Plan equivalence}: reflexivity and symmetry proved.
\end{enumerate}

The formalization comprises over 200 lines of Lean~4 code in
\texttt{Paper85\_AgentPlanningDescent.lean}, with all theorems
discharged without \texttt{sorry}.  The main summary theorem
\texttt{paper85\_summary} packages all results into a single
conjunction.


% =====================================================================
% Bibliography
% =====================================================================

\begin{thebibliography}{99}

\bibitem{paper00}
JuGeo Research Group.
\newblock Judgment geometry: A sheaf-theoretic foundation for
  software verification.
\newblock \emph{JuGeo Technical Report}, 2024.

\bibitem{paper01}
JuGeo Research Group.
\newblock Semantic sites for program analysis.
\newblock \emph{JuGeo Technical Report}, Paper~01, 2024.

\bibitem{paper03}
JuGeo Research Group.
\newblock Descent and obstructions in judgment geometry.
\newblock \emph{JuGeo Technical Report}, Paper~03, 2024.

\bibitem{paper04}
JuGeo Research Group.
\newblock Trust algebra: Ordered structures for verification trust.
\newblock \emph{JuGeo Technical Report}, Paper~04, 2024.

\bibitem{paper05}
JuGeo Research Group.
\newblock SMT dispatch for fragment-aware verification.
\newblock \emph{JuGeo Technical Report}, Paper~05, 2024.

\bibitem{paper07}
JuGeo Research Group.
\newblock Python effect tracking via judgment geometry.
\newblock \emph{JuGeo Technical Report}, Paper~07, 2024.

\bibitem{paper09}
JuGeo Research Group.
\newblock Proof-carrying Python: Lean~4 formalization of judgment
  geometry.
\newblock \emph{JuGeo Technical Report}, Paper~09, 2024.

\bibitem{paper24}
JuGeo Research Group.
\newblock Async effects in judgment geometry.
\newblock \emph{JuGeo Technical Report}, Paper~24, 2024.

\bibitem{paper37}
JuGeo Research Group.
\newblock Pack federation for distributed verification.
\newblock \emph{JuGeo Technical Report}, Paper~37, 2024.

\bibitem{paper42}
JuGeo Research Group.
\newblock Oracle federation: Multi-oracle judgment composition.
\newblock \emph{JuGeo Technical Report}, Paper~42, 2024.

\bibitem{paper43}
JuGeo Research Group.
\newblock Hypercovers and higher descent in judgment geometry.
\newblock \emph{JuGeo Technical Report}, Paper~43, 2024.

\bibitem{paper63}
JuGeo Research Group.
\newblock Planning code migrations using change-of-site functors.
\newblock \emph{JuGeo Technical Report}, Paper~63, 2024.

\bibitem{paper67}
JuGeo Research Group.
\newblock Code review automation via judgment ensembles.
\newblock \emph{JuGeo Technical Report}, Paper~67, 2024.

\bibitem{alfworld}
M.~Shridhar, X.~Yuan, M.-A.~C\^ot\'e, Y.~Bisk, A.~Trischler,
  and M.~Hausknecht.
\newblock {ALFWorld}: Aligning text and embodied environments for
  interactive learning.
\newblock In \emph{Proc.\ ICLR}, 2021.

\bibitem{webshop}
S.~Yao, H.~Chen, J.~Yang, and K.~Narasimhan.
\newblock {WebShop}: Towards scalable real-world web interaction with
  grounded language agents.
\newblock In \emph{Proc.\ NeurIPS}, 2022.

\bibitem{swebench}
C.~E. Jimenez, J.~Yang, A.~Wettig, S.~Yao, K.~Pei, O.~Press,
  and K.~Narasimhan.
\newblock {SWE}-bench: Can language models resolve real-world
  {GitHub} issues?
\newblock In \emph{Proc.\ ICLR}, 2024.

\bibitem{gpt4}
OpenAI.
\newblock {GPT-4} technical report.
\newblock \emph{arXiv preprint arXiv:2303.08774}, 2023.

\bibitem{cot}
J.~Wei, X.~Wang, D.~Schuurmans, M.~Bosma, B.~Ichter, F.~Xia,
  E.~Chi, Q.~Le, and D.~Zhou.
\newblock Chain-of-thought prompting elicits reasoning in large
  language models.
\newblock In \emph{Proc.\ NeurIPS}, 2022.

\bibitem{tot}
S.~Yao, D.~Yu, J.~Zhao, I.~Shafran, T.~L. Griffiths, Y.~Cao,
  and K.~Narasimhan.
\newblock Tree of thoughts: Deliberate problem solving with large
  language models.
\newblock In \emph{Proc.\ NeurIPS}, 2023.

\bibitem{planandsolve}
L.~Wang, W.~Xu, Y.~Lan, Z.~Hu, Y.~Lan, R.~K.-W. Lee, and
  E.-P. Lim.
\newblock Plan-and-solve prompting: Improving zero-shot
  chain-of-thought reasoning by large language models.
\newblock In \emph{Proc.\ ACL}, 2023.

\bibitem{react}
S.~Yao, J.~Zhao, D.~Yu, N.~Du, I.~Shafran, K.~Narasimhan,
  and Y.~Cao.
\newblock {ReAct}: Synergizing reasoning and acting in language
  models.
\newblock In \emph{Proc.\ ICLR}, 2023.

\bibitem{reflexion}
N.~Shinn, F.~Cassano, A.~Gopinath, K.~Narasimhan, and S.~Yao.
\newblock Reflexion: Language agents with verbal reinforcement
  learning.
\newblock In \emph{Proc.\ NeurIPS}, 2023.

\bibitem{howey2004val}
R.~Howey, D.~Long, and M.~Fox.
\newblock {VAL}: Automatic plan validation, continuous effects and
  mixed initiative planning using {PDDL}.
\newblock In \emph{Proc.\ ICTAI}, 2004.

\bibitem{pvp}
M.~Katz, S.~Sohrabi, and O.~Udrea.
\newblock Plan verification as satisfiability.
\newblock In \emph{Proc.\ ICAPS}, 2018.

\bibitem{robinson2014topological}
M.~Robinson.
\newblock \emph{Topological Signal Processing}.
\newblock Springer, 2014.

\bibitem{hansen2019opinion}
J.~Hansen and R.~Ghrist.
\newblock Opinion dynamics on discourse sheaves.
\newblock \emph{SIAM Journal on Applied Mathematics}, 81(5), 2021.

\bibitem{goguen1992sheaf}
J.~A. Goguen.
\newblock Sheaf semantics for concurrent interacting objects.
\newblock \emph{Mathematical Structures in Computer Science},
  2(2):159--191, 1992.

\bibitem{curry2014sheaves}
J.~Curry.
\newblock \emph{Sheaves, Cosheaves and Applications}.
\newblock PhD thesis, University of Pennsylvania, 2014.

\bibitem{ghallab2004automated}
M.~Ghallab, D.~Nau, and P.~Traverso.
\newblock \emph{Automated Planning: Theory and Practice}.
\newblock Morgan Kaufmann, 2004.

\bibitem{nau2004automated}
D.~S. Nau, M.~Ghallab, and P.~Traverso.
\newblock \emph{Automated Planning and Acting}.
\newblock Cambridge University Press, 2016.

\bibitem{gastin2001fast}
P.~Gastin and D.~Oddoux.
\newblock Fast {LTL} to {B\"uchi} automata translation.
\newblock In \emph{Proc.\ CAV}, 2001.

\bibitem{clarke2018handbook}
E.~M. Clarke, T.~A. Henzinger, H.~Veith, and R.~Bloem, editors.
\newblock \emph{Handbook of Model Checking}.
\newblock Springer, 2018.

\bibitem{dcop}
A.~Petcu and B.~Faltings.
\newblock A scalable method for multiagent constraint optimization.
\newblock In \emph{Proc.\ IJCAI}, 2005.

\bibitem{de2003plan}
M.~de~Weerdt, A.~ter~Mors, and C.~Witteveen.
\newblock Multi-agent planning: An introduction to planning and
  coordination.
\newblock In \emph{Handouts of the European Agent Systems Summer
  School}, 2003.

\end{thebibliography}

\end{document}
